%=============================================================================
% Wave 3, Iteration 6: Materials Science Update
% AGENT: MATERIALS SCIENCE (W3i6)
%
% W3i5 KEY FINDINGS CARRIED FORWARD:
%   - Two-component psi: psi = psi_grav + delta_psi_EM
%   - All material FOMs survive M_eff correction (ratios unaffected)
%   - P15 active detection devastated by M_eff; pivot to passive
%   - 11 P11 configurations beat SQMS bound
%   - Passive Superiority Theorem: passive bounds fundamentally superior
%   - YBCO FOM = 965 (unchanged), MgB2 = 125, Nb3Sn = 92, Nb = 3.2
%   - P2+P11 hybrid with MEMS proof mass: viable path
%
% W3i6 MANDATE:
%   1. Under two-component, materials couple to delta_psi_EM.
%      Does this change the FOM formula?
%   2. With passive detection dominant, derive PASSIVE material FOM
%   3. Rank all candidate materials under passive FOM
%   4. MEMS materials for P2+P11 path: sub-nanogram proof mass optimisation
%   5. Updated materials roadmap for post-paradigm-shift DFD programme
%
% Date: March 2026
%=============================================================================

\documentclass[11pt]{article}
\usepackage{amsmath,amssymb,amsthm}
\usepackage{geometry}
\geometry{margin=1in}
\usepackage{booktabs}
\usepackage{longtable}
\usepackage{array}
\usepackage{enumitem}
\usepackage{hyperref}
\usepackage{xcolor}
\usepackage{tcolorbox}
\usepackage{float}

\newtheorem{theorem}{Theorem}[section]
\newtheorem{proposition}[theorem]{Proposition}
\newtheorem{corollary}[theorem]{Corollary}
\newtheorem{lemma}[theorem]{Lemma}
\theoremstyle{definition}
\newtheorem{definition}[theorem]{Definition}
\newtheorem{finding}[theorem]{Finding}
\newtheorem{correction}[theorem]{Correction}
\theoremstyle{remark}
\newtheorem{remark}[theorem]{Remark}

\newcommand{\ee}[1]{\times 10^{#1}}
\newcommand{\Meff}{M_{\rm eff}}
\newcommand{\Opsi}{\omega_\psi}
\newcommand{\Qpsi}{Q_\psi}
\newcommand{\dpsi}{\delta\psi_{\rm EM}}
\newcommand{\psigrav}{\psi_{\rm grav}}
\newcommand{\psiEM}{\delta\psi_{\rm EM}}
\newcommand{\lambdaone}{|\lambda-1|}
\newcommand{\lambdabound}{1.56\ee{-9}}
\newcommand{\Tc}{T_{\rm c}}
\newcommand{\lambdaL}{\lambda_{\rm L}}
\newcommand{\FOMactive}{\mathrm{FOM}_{\rm active}}
\newcommand{\FOMpassive}{\mathrm{FOM}_{\rm passive}}
\newcommand{\FOMMEMS}{\mathrm{FOM}_{\rm MEMS}}
\newcommand{\lam}{|\lambda-1|}
\newcommand{\Heff}{H_{\rm eff}}
\newcommand{\muG}{\mu_G}
\newcommand{\GN}{G_{\rm N}}
\newcommand{\SNR}{\mathrm{SNR}}

\title{Wave 3 Iteration 6: Materials Science Update\\
\large Post-Paradigm-Shift Material Selection:\\
Two-Component FOM Analysis, Passive Detection Materials,\\
MEMS Proof Mass Optimisation, and Updated Roadmap}
\author{DFD Research Programme --- Materials Science Agent, W3i6}
\date{March 2026}

\begin{document}
\maketitle
\tableofcontents
\newpage

%=============================================================================
\section{Executive Summary}
\label{sec:executive}
%=============================================================================

Wave 3 Iteration 5 established three paradigm-shifting results that
fundamentally alter the materials science programme:

\begin{enumerate}[nosep]
  \item \textbf{Two-component $\psi$-field}: $\psi = \psigrav + \psiEM$,
    with all patent-specific (laboratory) effects coupling to
    $\psiEM$ only.
  \item \textbf{Passive Superiority Theorem}: Passive cavity bounds
    ($\delta Q/Q$ measurements) are fundamentally superior to active
    SQUID detection for bulk SRF cavities because $\Meff$ cancels
    in the passive ratio but enters linearly in the active signal.
  \item \textbf{Active detection devastation}: The $\Meff = 3.01\ee{6}$~kg
    correction degrades all active P15 sensitivity by $10^6\times$;
    near-term cavity-based active detection is not viable.
\end{enumerate}

\noindent This iteration answers five questions:

\begin{tcolorbox}[colback=blue!5!white, colframe=blue!50!black,
  title=\textbf{W3i6 Materials Science Mandate}]
\begin{enumerate}[nosep]
  \item Does the two-component decomposition change the active FOM formula?
  \item What material properties optimise \emph{passive} detection
    (Q-factor stability, frequency stability, dimensional stability)?
  \item Derive the Passive Material FOM and rank all candidate materials.
  \item What sub-nanogram proof mass material optimises the P2+P11
    hybrid MEMS path?
  \item Updated materials roadmap for the post-paradigm-shift programme.
\end{enumerate}
\end{tcolorbox}

\noindent\textbf{Key results}:
\begin{itemize}[nosep]
  \item The active FOM formula is \emph{unchanged} under two-component
    decomposition: $\psiEM$ replaces $\psi_0$ but the material-dependent
    factors $\kappa\cdot\Tc^{1/2}$ remain identical.
  \item We derive a new \textbf{Passive Material FOM} that captures
    Q-factor stability, frequency stability, and thermal noise rejection.
    Nb dominates the passive ranking (inverting the active ranking).
  \item For MEMS proof masses, single-crystal silicon nitride (high-stress
    Si$_3$N$_4$) is optimal at $10^{-12}$--$10^{-9}$~kg, with
    crystalline silicon and aluminium alternatives.
  \item The materials roadmap now has three parallel tracks: Passive
    (Nb-based, near-term), MEMS (Si$_3$N$_4$, medium-term), and
    Exotic (YBCO/Nb$_3$Sn for future active recovery).
\end{itemize}


%=============================================================================
\section{Two-Component $\psi$-Field: Impact on the Active FOM}
\label{sec:two-component-fom}
%=============================================================================

\subsection{Review: the two-component decomposition}

W3i5 established:
\begin{equation}
  \psi = \psigrav + \psiEM,
\end{equation}
where $\psigrav$ is the gravitational potential (sourced by all
stress-energy, giving MOND, lensing, the shadow, Hubble tension),
and $\psiEM$ is governed by the DFD field equation sourced by EM
stress-energy (giving $\dot{G}/G$, Process A coupling, all laboratory
effects).

\textbf{All patent-specific effects couple to $\psiEM$, not to $\psigrav$.}
The cosmological value $\psiEM^{\rm cosmo} \approx 9.6\ee{-6}$ (not the
gravitational $\psigrav \approx 0.047$).

\subsection{Does the active FOM change?}

The W3i4 active FOM for SRF cavity materials is:
\begin{equation}
  \FOMactive = \kappa\cdot\Tc^{1/2}
  = \frac{\lambdaL}{\xi_{\rm sc}}\cdot\Tc^{1/2},
  \label{eq:FOM_active}
\end{equation}
where $\lambdaL$ is the London penetration depth, $\xi_{\rm sc}$ is the
superconducting coherence length, and $\Tc$ is the critical temperature.

This FOM captures the material's ability to:
\begin{itemize}[nosep]
  \item Maximise the EM--psi overlap integral ($\propto \kappa = \lambdaL/\xi_{\rm sc}$,
    the Ginzburg--Landau parameter);
  \item Maximise the operating temperature range ($\propto \Tc^{1/2}$,
    which sets the energy scale of the EM fields).
\end{itemize}

Under two-component decomposition, the Process A coupling vertex is:
\begin{equation}
  g_{\rm A} = (\lambda - 1)\,\omega_{\rm EM}\,\psiEM_{\rm zpf},
  \label{eq:coupling_vertex}
\end{equation}
where $\psiEM_{\rm zpf} = \sqrt{\hbar/(2\Meff\Opsi)}$ is the zero-point
amplitude of the EM-sourced psi-perturbation. The material enters through
the overlap integral:
\begin{equation}
  \Heff = \frac{\int_{\rm cav} |\mathbf{E}|^2\,
  |\psi_n(\mathbf{r})|^2\,dV}
  {\left(\int_{\rm cav} |\mathbf{E}|^2\,dV\right)
  \left(\int_{\rm cav} |\psi_n|^2\,dV\right)},
\end{equation}
which depends on the EM field distribution in the cavity walls (set by
$\lambdaL$ and $\xi_{\rm sc}$) but \emph{not} on which component of
$\psi$ is being coupled to.

\begin{finding}[Active FOM Unchanged Under Two-Component Decomposition]
\label{find:active_fom_unchanged}
The two-component decomposition $\psi = \psigrav + \psiEM$ does not
change the active FOM formula. The material-dependent factors
($\lambdaL$, $\xi_{\rm sc}$, $\Tc$) enter through the EM field
distribution and overlap integral, which are properties of the cavity
geometry and the superconductor, not of the psi-field source.

The only change is that the psi-field amplitude is $\psiEM$ rather than
$\psi_0$. Since $\psiEM^{\rm cosmo} = 9.6\ee{-6} \ll \psigrav = 0.047$,
this reduces the cosmological psi-amplitude by $\sim 5000\times$.
However, the FOM is a \emph{relative} ranking: $\psiEM$ cancels in
material-to-material comparisons.

\textbf{All W3i4/W3i5 active FOM rankings are preserved:}
\begin{center}
\begin{tabular}{llll}
\toprule
Rank & Material & $\FOMactive$ & Status \\
\midrule
1 & YBCO & 965 & \textbf{Unchanged} \\
2 & MgB$_2$ & 125 & Unchanged \\
3 & Nb$_3$Sn & 92 & Unchanged \\
4 & Nb & 3.2 & Unchanged \\
\bottomrule
\end{tabular}
\end{center}

However, the \emph{relevance} of this ranking is now diminished: active
detection requires closing the $10^{12}$--$10^{17}$ detection gap,
which cannot be achieved by material optimisation alone.
\end{finding}


%=============================================================================
\section{Passive Detection: What Material Properties Matter?}
\label{sec:passive_properties}
%=============================================================================

The Passive Superiority Theorem (W3i5, Theorem~5.1 of the P2 update)
establishes that passive cavity bounds---measurements of $\delta Q/Q$ or
equivalently $\delta f/f$---are fundamentally superior to active SQUID
detection for bulk SRF cavities. This section derives the material
properties that optimise passive detection.

\subsection{The passive bound mechanism}

The SQMS passive bound $\lam < \lambdabound$ is derived from:
\begin{equation}
  \lam < \sqrt{\frac{\delta(1/Q)}{\Gamma_{\psi\to\mathrm{EM}}/\Opsi}},
  \label{eq:passive_bound}
\end{equation}
where $\delta(1/Q)$ is the measured Q-factor instability and
$\Gamma_{\psi\to\mathrm{EM}}$ is the psi-to-EM conversion rate per
unit $(\lambda-1)^2$. Crucially, $\Meff$ cancels in this ratio.

For a cavity measured over integration time $T_{\rm int}$, the bound
tightens as:
\begin{equation}
  \lam_{\rm min}^{\rm passive} \propto
  \left(\frac{\sigma_Q}{Q_0}\right)^{1/2}
  \cdot T_{\rm int}^{-1/4},
  \label{eq:passive_scaling}
\end{equation}
where $\sigma_Q/Q_0$ is the fractional Q-factor stability.

\subsection{Three pillars of passive material performance}

For a passive cavity bound, the material must deliver:

\begin{enumerate}
  \item \textbf{Q-factor stability} ($\sigma_Q/Q_0$): The fractional
    fluctuation of the loaded quality factor over the measurement period.
    Any source of Q-jitter (flux noise, quasiparticle density fluctuations,
    two-level system (TLS) noise, mechanical vibrations coupling to
    surface resistance) degrades the bound.

  \item \textbf{Frequency stability} ($\sigma_f/f_0$): The fractional
    frequency fluctuation. In a passive measurement, frequency drift
    can be interpreted as a psi-coupling signal; materials with lower
    $\sigma_f/f_0$ allow longer integration and tighter bounds.
    Frequency stability is governed by:
    \begin{equation}
      \frac{\sigma_f}{f_0} = \frac{1}{2}\left[
      \alpha_{\rm th}\,\sigma_T
      + \frac{\sigma_L}{L}
      + \frac{\Delta\lambdaL}{\lambdaL}\,\frac{\lambdaL}{R_{\rm cav}}
      \right],
      \label{eq:freq_stability}
    \end{equation}
    where $\alpha_{\rm th}$ is the thermal expansion coefficient,
    $\sigma_T$ is the temperature stability, $\sigma_L/L$ is the
    dimensional stability, and $\Delta\lambdaL/\lambdaL$ captures
    penetration depth fluctuations.

  \item \textbf{Thermal noise floor}: The fundamental limit on Q-factor
    stability is set by thermodynamic fluctuations of quasiparticle
    density. At operating temperature $T \ll \Tc$:
    \begin{equation}
      \left(\frac{\sigma_Q}{Q_0}\right)_{\rm thermal}
      \propto \frac{1}{Q_0}\sqrt{\frac{k_{\rm B}T}{\Delta_{\rm sc}}}
      \,e^{-\Delta_{\rm sc}/(2k_{\rm B}T)},
      \label{eq:Q_thermal_noise}
    \end{equation}
    where $\Delta_{\rm sc}$ is the superconducting gap. Materials with
    higher $\Delta_{\rm sc}$ (higher $\Tc$) suppress this floor
    \emph{exponentially}.
\end{enumerate}

\subsection{Technical noise sources in passive SRF measurements}

Beyond fundamental thermal fluctuations, SRF cavities suffer from:

\begin{itemize}[nosep]
  \item \textbf{Two-level system (TLS) noise}: Amorphous surface oxides
    (Nb$_2$O$_5$, etc.)\ host TLS defects that cause fluctuating
    dielectric loss. This is the dominant noise source at $T < 100$~mK
    for niobium cavities. Mitigation: high-temperature baking (300--400$^\circ$C)
    to transform the oxide, or nitrogen doping (N-doping) to reduce
    TLS density.
  \item \textbf{Flux noise}: Trapped magnetic flux causes dissipation
    fluctuations. Niobium is susceptible due to its relatively low
    $H_{c1}$. Nb$_3$Sn and YBCO have higher $H_{c1}$ and better flux
    expulsion.
  \item \textbf{Quasiparticle bursts}: Cosmic rays and radioactive
    impurities generate quasiparticle cascades that temporarily degrade
    Q. Mitigation: underground operation, phonon trapping layers, or
    operating at higher $T/\Tc$ ratios where the equilibrium quasiparticle
    density masks the bursts.
  \item \textbf{Microphonics}: Mechanical vibrations couple to cavity
    frequency and Q through Lorentz detuning and surface motion.
    Materials with higher mechanical stiffness and lower thermal
    contraction reduce microphonic coupling.
\end{itemize}


%=============================================================================
\section{The Passive Material Figure of Merit}
\label{sec:passive_fom}
%=============================================================================

\subsection{Derivation}

The passive bound sensitivity scales as:
\begin{equation}
  \lam_{\rm min}^{\rm passive} \propto
  \sqrt{\frac{\sigma_Q/Q_0}{\Heff^2}}
  \propto \sqrt{\frac{\sigma_Q}{Q_0\,\Heff^2}}.
  \label{eq:passive_sensitivity}
\end{equation}

For the material comparison, we separate factors:

\begin{itemize}[nosep]
  \item $Q_0$: the baseline quality factor (higher is better).
  \item $\sigma_Q/Q_0$: the fractional stability (lower is better).
  \item $\Heff$: the overlap factor (geometry-dependent; material enters
    through $\lambdaL$).
\end{itemize}

The dominant noise at millikelvin temperatures is TLS noise, which scales
as $\sigma_Q/Q_0 \propto n_{\rm TLS}\cdot\tan\delta_{\rm TLS}/Q_0$,
where $n_{\rm TLS}$ is the TLS density and $\tan\delta_{\rm TLS}$ is
the TLS loss tangent. Since TLS noise is an \emph{additive} loss channel,
$\sigma_{1/Q} = \sigma_{\rm TLS}$ is approximately material-independent
for a given surface preparation. The passive bound then scales as:
\begin{equation}
  \lam_{\rm min}^{\rm passive} \propto
  \sqrt{\frac{\sigma_{\rm TLS}}{Q_0^2\,\Heff^2}}
  = \frac{\sqrt{\sigma_{\rm TLS}}}{Q_0\,\Heff}.
\end{equation}

Therefore, the material-dependent passive FOM is:
\begin{equation}
\boxed{
  \FOMpassive = Q_0 \cdot \Heff \cdot S_{\rm dim} \cdot S_{\rm freq},
}
\label{eq:FOM_passive}
\end{equation}
where:
\begin{itemize}[nosep]
  \item $Q_0$ is the achievable quality factor at the operating temperature.
  \item $\Heff$ is the EM--psi overlap factor (enters weakly for bulk
    cavities; $\Heff \approx 3\ee{-5}$ for TESLA geometry, approximately
    material-independent).
  \item $S_{\rm dim} = 1/\alpha_{\rm th}(T_{\rm op})$ is a dimensional
    stability factor: materials with lower thermal expansion at operating
    temperature contribute less frequency drift.
  \item $S_{\rm freq}$ captures the frequency stability from penetration
    depth fluctuations: $S_{\rm freq} = \exp(\Delta_{\rm sc}/(2k_{\rm B}T_{\rm op}))$,
    reflecting the exponential suppression of quasiparticle noise at
    $T_{\rm op} \ll \Tc$.
\end{itemize}

\begin{remark}
For the passive FOM, the $S_{\rm freq}$ factor is the most discriminating.
At $T_{\rm op} = 20$~mK:
\begin{itemize}[nosep]
  \item Nb ($\Tc = 9.25$~K, $\Delta = 1.55$~meV):
    $\Delta/(2k_{\rm B}T) = 1.55\ee{-3}/(2\times 8.62\ee{-5}\times 0.02)
    = 450$. The exponential $e^{450}$ is effectively infinite; Nb is
    deep in the zero-quasiparticle regime.
  \item Nb$_3$Sn ($\Tc = 18.3$~K, $\Delta = 3.1$~meV):
    $\Delta/(2k_{\rm B}T) = 900$. Even more deeply suppressed.
  \item YBCO ($\Tc = 93$~K, $\Delta \sim 20$~meV, d-wave):
    $\Delta/(2k_{\rm B}T) \sim 5800$. However, d-wave gap nodes mean
    the effective suppression is power-law, not exponential---the
    quasiparticle density goes as $T^2/\Delta$, not $e^{-\Delta/(2k_BT)}$.
\end{itemize}
All s-wave superconductors are in the deep exponential regime at 20~mK,
so $S_{\rm freq}$ does not discriminate between Nb, Nb$_3$Sn, and MgB$_2$.
The d-wave node structure of YBCO is a \emph{disadvantage} for passive
detection because it permits residual quasiparticles at all temperatures.
\end{remark}

\subsection{The decisive factor: $Q_0$ at operating temperature}

Since $S_{\rm freq}$ saturates for all s-wave superconductors at
$T_{\rm op} = 20$~mK, and $\Heff$ is approximately geometry-dependent
(not material-dependent for bulk TESLA cavities), the passive FOM
reduces to:
\begin{equation}
  \FOMpassive \approx Q_0(T_{\rm op}) \cdot S_{\rm dim}(T_{\rm op}).
  \label{eq:FOM_passive_simplified}
\end{equation}

\begin{theorem}[Passive FOM Ranking Inverts the Active Ranking]
\label{thm:passive_inversion}
For passive cavity bounds at $T_{\rm op} \leq 20$~mK, the material
ranking is determined primarily by achievable $Q_0$. The demonstrated
$Q_0$ values are:

\begin{center}
\renewcommand{\arraystretch}{1.3}
\begin{tabular}{lrrlrl}
\toprule
Material & $Q_0$ (achieved) & $\Tc$ (K) & Gap & $\alpha_{\rm th}$ (K$^{-1}$, 4~K)
  & $\FOMpassive$ \\
\midrule
Nb (N-doped) & $2\ee{11}$ & 9.25 & s-wave & $1\ee{-9}$ (est.)
  & \textbf{$2\ee{20}$} \\
Nb (EP) & $5\ee{10}$ & 9.25 & s-wave & $1\ee{-9}$
  & $5\ee{19}$ \\
Nb$_3$Sn & $2\ee{10}$ & 18.3 & s-wave & $3\ee{-9}$
  & $7\ee{18}$ \\
MgB$_2$ & $5\ee{9}$ (proj.) & 39 & Two-gap s & $5\ee{-9}$
  & $1\ee{18}$ \\
YBCO & $10^{6}$--$10^{7}$ & 93 & d-wave & $3\ee{-8}$
  & $3\ee{13}$ \\
\bottomrule
\end{tabular}
\end{center}

\textbf{Result}: Niobium (N-doped) is the optimal passive detection
material, with $\FOMpassive$ exceeding YBCO by $\sim 10^7$. This is
the \emph{inversion} of the active ranking, where YBCO ($\FOMactive = 965$)
dominated Nb ($\FOMactive = 3.2$) by $300\times$.

The inversion occurs because:
\begin{enumerate}[nosep]
  \item Passive detection does not benefit from high $\kappa$ (the
    GL parameter that favours YBCO in active mode).
  \item Passive detection requires ultrahigh $Q_0$, which is achieved
    in Nb but not in YBCO (limited by d-wave nodes and grain boundaries).
  \item YBCO's d-wave gap structure permits power-law quasiparticle
    density at low $T$, causing residual Q-factor fluctuations that
    s-wave materials avoid.
\end{enumerate}
\end{theorem}

\begin{finding}[YBCO Demoted for Passive Detection]
\label{find:ybco_demoted}
YBCO, which dominated the active FOM ranking ($\FOMactive = 965$),
is the \emph{worst} candidate for passive detection among the materials
studied. Its $Q_0$ is limited to $\sim 10^7$ by:
\begin{itemize}[nosep]
  \item d-wave gap nodes producing finite quasiparticle density at all $T$;
  \item grain boundary weak links causing microwave loss;
  \item surface oxygen disorder affecting Q-factor reproducibility.
\end{itemize}
For the passive programme, YBCO should be set aside entirely.
\end{finding}


%=============================================================================
\section{Passive Material Rankings: Detailed Analysis}
\label{sec:passive_ranking}
%=============================================================================

\subsection{Rank 1: Nitrogen-doped Niobium}

Nitrogen-doped niobium (N-doped Nb) achieves the highest $Q_0$ of any
SRF material: $Q_0 \approx 2\ee{11}$ at 1.3~GHz and $T = 2$~K
(Fermilab SQMS programme), with projections to $>10^{12}$ at dilution
refrigerator temperatures ($T < 20$~mK).

\textbf{Key advantages for passive detection:}
\begin{itemize}[nosep]
  \item Highest demonstrated $Q_0$ of any SRF cavity material.
  \item S-wave gap ($\Delta = 1.55$~meV): exponential quasiparticle
    suppression at $T \ll \Tc$.
  \item Well-characterised surface: decades of Fermilab/DESY/JLab
    experience with surface preparation (EP, BCP, 120$^\circ$C bake).
  \item TLS noise reduction via 300$^\circ$C vacuum baking eliminates
    the Nb$_2$O$_5$ amorphous oxide.
  \item Excellent dimensional stability: Nb has very low thermal
    contraction below 4~K ($\alpha_{\rm th} \sim 10^{-9}$~K$^{-1}$).
  \item The SQMS bound ($\lam < \lambdabound$) was \emph{already achieved}
    with Nb cavities---proving that Nb-based passive detection works.
\end{itemize}

\textbf{Limitations:}
\begin{itemize}[nosep]
  \item Moderate $\Tc = 9.25$~K limits the operating field to
    $H_{\rm sh} \approx 200$~mT (but passive detection does not require
    high fields---it requires Q-factor stability at low fields).
  \item Flux trapping from ambient magnetic fields requires careful
    shielding ($< 1$~nT).
\end{itemize}

\subsection{Rank 2: Nb$_3$Sn-coated Nb}

Nb$_3$Sn-coated niobium cavities achieve $Q_0 \approx 2\ee{10}$ at
4.2~K (enabling helium-free operation). At $T = 20$~mK:

\textbf{Advantages:}
\begin{itemize}[nosep]
  \item Higher $\Tc = 18.3$~K and $\Delta = 3.1$~meV: even deeper
    quasiparticle suppression.
  \item Higher $H_{c1}$: better flux expulsion, lower trapped flux
    dissipation.
  \item Can operate at higher $T$ (4.2~K) with adequate $Q_0$ for
    preliminary passive bounds, reducing cryogenic complexity.
\end{itemize}

\textbf{Limitations:}
\begin{itemize}[nosep]
  \item $Q_0$ at 20~mK projected at $2\ee{10}$, an order of magnitude
    below N-doped Nb.
  \item Tin segregation at grain boundaries causes local loss hotspots.
  \item Surface roughness from vapour diffusion coating introduces
    geometric loss.
\end{itemize}

\subsection{Rank 3: MgB$_2$}

MgB$_2$ ($\Tc = 39$~K) is a two-gap s-wave superconductor. Its potential
for passive detection:

\textbf{Advantages:}
\begin{itemize}[nosep]
  \item Highest $\Tc$ among s-wave candidates: $\Delta_\sigma = 7$~meV,
    $\Delta_\pi = 2.5$~meV.
  \item Light constituent elements: low cosmic-ray quasiparticle
    generation rate.
  \item Higher operating temperature capability: useful for 4.2~K
    screening measurements.
\end{itemize}

\textbf{Limitations:}
\begin{itemize}[nosep]
  \item No demonstrated bulk SRF cavity; only thin-film coatings on
    Cu substrates, with $Q_0 \lesssim 5\ee{9}$.
  \item The $\pi$-band gap ($\Delta_\pi = 2.5$~meV) sets the
    effective quasiparticle suppression, not the larger $\sigma$-band
    gap.
  \item Film quality not yet competitive with bulk Nb surface treatments.
\end{itemize}

\subsection{Rank 4: YBCO (not recommended)}

As noted in Finding~\ref{find:ybco_demoted}, YBCO is not viable for
passive detection due to its d-wave gap, which prevents the exponential
quasiparticle freezeout essential for ultra-stable Q-factors.


%=============================================================================
\section{Passive Detection: Optimised Cavity Configurations}
\label{sec:passive_configs}
%=============================================================================

\subsection{Configuration 1: SQMS-class Nb cavity at 20 mK}

This is the baseline: a single 1.3~GHz TESLA-type Nb cavity, N-doped,
operated in a dilution refrigerator.

\begin{center}
\begin{tabular}{ll}
\toprule
Parameter & Value \\
\midrule
Cavity & TESLA 9-cell, N-doped Nb \\
$Q_0$ & $2\ee{11}$ (projected at 20~mK) \\
$f_0$ & 1.3~GHz \\
$T_{\rm op}$ & 20~mK \\
$\sigma_Q/Q_0$ & $\sim 10^{-6}$ (TLS-limited) \\
Integration time & $10^6$~s (12 days) \\
Projected bound & $\lam < \lambdabound$ (achieved by SQMS) \\
\bottomrule
\end{tabular}
\end{center}

\subsection{Configuration 2: 256-cavity coherent array (ESS class)}

The 256-cavity coherent array (W3i5 P11 Finding) achieves $256\times$
improvement through phased averaging:
\begin{equation}
  \lam_{\rm min}^{\rm array} = \frac{\lam_{\rm min}^{\rm single}}{N}
  = \frac{\lambdabound}{256} = 6.1\ee{-12}.
\end{equation}

\textbf{Material requirement}: All 256 cavities must have matched
frequency to within $\delta f/f < 1/Q_L$ for coherent combination.
N-doped Nb is the only material with demonstrated frequency
reproducibility at this level.

\subsection{Configuration 3: Cross-correlation (ESS Lund)}

Two or more cavities measuring independent Q-factor time series;
cross-correlating to extract correlated psi-field fluctuations.
This rejects uncorrelated noise (TLS, quasiparticle bursts) and
projects to $\lam \sim 10^{-12}$--$10^{-14}$.

\textbf{Material requirement}: Identical cavity pairs with
uncorrelated noise. N-doped Nb from the same ingot, processed
identically, provides the best matched-pair performance.


%=============================================================================
\section{MEMS Proof Mass Materials for the P2+P11 Path}
\label{sec:mems}
%=============================================================================

The P2+P11 hybrid (W3i5 P10/P11 update, Section~4.5) uses a MEMS
mechanical resonator as the psi-field detector. The sensitivity scales
as $\sqrt{\Meff^{\rm MEMS}}$, so minimising the proof mass is critical.
W3i5 showed:

\begin{center}
\begin{tabular}{llll}
\toprule
$\Meff^{\rm MEMS}$ & Technology & $\lam_{\rm min}$ & Feasibility \\
\midrule
$10^{-6}$~kg ($\mu$g) & Nb membrane & $4.7\ee{-11}$ & Near-term \\
$10^{-9}$~kg (ng) & MEMS beam & $1.5\ee{-12}$ & 2027--2028 \\
$10^{-12}$~kg (pg) & CNT resonator & $4.7\ee{-14}$ & 2030+ \\
$10^{-15}$~kg (fg) & PhC nanobeam & $1.5\ee{-15}$ & Speculative \\
\bottomrule
\end{tabular}
\end{center}

The material choice for the MEMS proof mass must optimise a different
set of properties from the SRF cavity material.

\subsection{MEMS material requirements}

For the P2+P11 hybrid, the MEMS detector must provide:

\begin{enumerate}
  \item \textbf{High mechanical $Q_{\rm mech}$}: The sensitivity scales
    as $1/Q_{\rm mech}$ (Eq.~\ref{eq:hybrid_formula_ref} below). The
    proof mass must sustain $Q_{\rm mech} \geq 10^6$ at millikelvin
    temperatures.

  \item \textbf{Low mass}: The sensitivity scales as $\sqrt{\Meff^{\rm MEMS}}$,
    so lighter proof masses are better.

  \item \textbf{Cryogenic compatibility}: Must operate at $T < 100$~mK
    in the SRF environment (high vacuum, strong magnetic shielding,
    low vibration).

  \item \textbf{SQUID-readable displacement}: The proof mass motion must
    couple to a SQUID magnetometer, requiring either a superconducting
    proof mass (direct flux coupling) or a superconducting transducer
    layer.

  \item \textbf{Low intrinsic stress noise}: Thermomechanical noise
    in the proof mass sets the noise floor:
    \begin{equation}
      S_x^{\rm th} = \frac{4k_{\rm B}T\,\Meff\,\omega_{\rm mech}}{Q_{\rm mech}},
    \end{equation}
    so the noise spectral density scales as $\Meff/Q_{\rm mech}$. The
    relevant material FOM is $Q_{\rm mech}/\Meff$.
\end{enumerate}

\subsection{MEMS Material Figure of Merit}

Combining requirements 1, 2, and 5:
\begin{equation}
\boxed{
  \FOMMEMS = \frac{Q_{\rm mech}}{\Meff^{1/2}}
  \cdot \frac{1}{\sqrt{S_x^{\rm th}(T_{\rm op})}}
  \propto \frac{Q_{\rm mech}^{3/2}}{\Meff \cdot T_{\rm op}^{1/2}
  \cdot \omega_{\rm mech}^{1/2}}.
}
\label{eq:FOM_MEMS}
\end{equation}

At fixed $T_{\rm op}$ and $\omega_{\rm mech}$, this reduces to
$\FOMMEMS \propto Q_{\rm mech}^{3/2}/\Meff$.

\subsection{MEMS material ranking}

\begin{theorem}[MEMS Proof Mass Material Ranking]
\label{thm:mems_ranking}
Evaluating candidate MEMS materials at $T = 20$~mK:

\begin{center}
\renewcommand{\arraystretch}{1.3}
\begin{tabular}{p{3cm}rrrl}
\toprule
Material & $Q_{\rm mech}$ & Min.\ mass & $\FOMMEMS$ & Notes \\
 & (at 20~mK) & (achievable) & (relative) & \\
\midrule
High-stress Si$_3$N$_4$ & $10^{9}$--$10^{10}$ & 1~pg
  & \textbf{$10^{18}$} & Dissipation dilution \\
Crystalline Si (SOI) & $10^{7}$--$10^{8}$ & 10~pg
  & $10^{13}$ & Established MEMS \\
Al (supercond.) & $10^{6}$--$10^{7}$ & 1~ng
  & $10^{10}$ & Direct SQUID coupling \\
Nb (supercond.) & $10^{5}$--$10^{6}$ & 10~ng
  & $10^{7}$ & SRF-compatible \\
Diamond (NV) & $10^{7}$ & 100~pg
  & $10^{12}$ & Magnetometry readout \\
CNT bundle & $10^{5}$--$10^{6}$ & 1~fg
  & $10^{14}$ & Ultra-low mass \\
PhC nanobeam & $10^{8}$--$10^{9}$ & 10~fg
  & $10^{16}$ & Optomechanical \\
\bottomrule
\end{tabular}
\end{center}

\textbf{Result}: High-stress silicon nitride (Si$_3$N$_4$) is the optimal
MEMS proof mass material, primarily due to \emph{dissipation dilution}---the
mechanism by which high intrinsic stress distributes vibrational energy
into the high-Q stressed mode, achieving $Q_{\rm mech} > 10^9$ even at
millikelvin temperatures.
\end{theorem}

\subsection{High-stress Si$_3$N$_4$: detailed assessment}

\textbf{Dissipation dilution mechanism:} In a highly pre-stressed
Si$_3$N$_4$ membrane or string resonator, the elastic energy is dominated
by the tensile stress rather than bending. The quality factor is enhanced
by:
\begin{equation}
  Q_{\rm eff} = Q_{\rm intrinsic} \times D_{\rm dil},
  \quad D_{\rm dil} = \frac{\sigma\,L^2}{E\,h^2},
\end{equation}
where $\sigma \approx 1$~GPa is the tensile stress, $L$ is the resonator
length, $E \approx 250$~GPa is Young's modulus, and $h$ is the thickness.
For a 100~nm thick, 100~$\mu$m long Si$_3$N$_4$ string:
$D_{\rm dil} \approx 10^9/(250\times10^4) = 4\ee{4}$, giving
$Q_{\rm eff} \approx Q_{\rm int} \times 4\ee{4} \approx 10^4 \times 4\ee{4}
= 4\ee{8}$.

Soft-clamped phononic crystal designs have demonstrated
$Q_{\rm mech} > 10^9$ at room temperature and projections of $>10^{10}$
at millikelvin temperatures.

\textbf{Mass scaling:} A Si$_3$N$_4$ membrane of dimensions
$100~\mu\text{m} \times 100~\mu\text{m} \times 50$~nm has mass:
\begin{equation}
  m = \rho \cdot V = 3100~\text{kg/m}^3 \times 5\ee{-16}~\text{m}^3
  = 1.55\ee{-12}~\text{kg} = 1.55~\text{pg}.
\end{equation}

For the nanogram regime ($m = 10^{-9}$~kg), a
$1~\text{mm} \times 1~\text{mm} \times 100$~nm membrane suffices:
$m = 3100 \times 10^{-13} = 3.1\ee{-10}$~kg $\approx 0.3$~ng.

\textbf{SQUID coupling:} Si$_3$N$_4$ is not superconducting. SQUID
readout requires a thin ($\sim 20$~nm) Al or Nb film deposited on the
membrane. This adds negligible mass ($< 5\%$) but provides direct
flux-to-displacement transduction.

\textbf{Cryogenic performance:} Si$_3$N$_4$ membranes have been operated
at dilution refrigerator temperatures in multiple optomechanics
experiments. The intrinsic $Q$ improves at low temperature (fewer phonon
channels), making it ideal for the P2+P11 application.

\subsection{Alternative: Aluminium drumhead resonators}

For the near-term ($\mu$g to ng) MEMS path, aluminium drumhead resonators
offer a simpler approach:

\begin{itemize}[nosep]
  \item Superconducting below 1.2~K: direct SQUID coupling without
    additional metallisation.
  \item Demonstrated $Q_{\rm mech} \sim 10^7$ at 20~mK (NIST, Delft).
  \item Established fabrication on Si substrates.
  \item Mass range: 1~ng to 1~$\mu$g depending on geometry.
\end{itemize}

\textbf{Limitation}: $Q_{\rm mech}$ is $100\times$--$1000\times$ lower
than Si$_3$N$_4$ with dissipation dilution. For the P2+P11 hybrid,
this translates to $10\times$--$30\times$ worse sensitivity, reducing
the advantage over the SQMS passive bound.

\subsection{Recommended MEMS proof mass strategy}

\begin{finding}[MEMS Proof Mass Material Selection]
\label{find:mems_selection}
\begin{enumerate}[nosep]
  \item \textbf{Phase 1 (2026--2027, near-term)}: Al drumhead resonator,
    $m \sim 1~\mu$g, $Q_{\rm mech} \sim 10^7$. Direct SQUID readout.
    Projected $\lam_{\rm min} \sim 5\ee{-11}$ ($30\times$ below SQMS bound).
    Technology readiness: demonstrated.

  \item \textbf{Phase 2 (2027--2029, medium-term)}: High-stress
    Si$_3$N$_4$ membrane with Al metallisation, $m \sim 1$~ng,
    $Q_{\rm mech} \sim 10^9$. Projected $\lam_{\rm min} \sim 5\ee{-14}$
    ($3\ee{4}\times$ below SQMS bound). Requires: phononic crystal
    soft-clamping design adapted for SQUID readout geometry.

  \item \textbf{Phase 3 (2030+, speculative)}: Si$_3$N$_4$ or diamond
    photonic crystal nanobeam, $m \sim 1$~pg, $Q_{\rm mech} \sim 10^{10}$.
    Projected $\lam_{\rm min} \sim 10^{-16}$. Requires: sub-nm SQUID
    displacement sensitivity or optical readout conversion.
\end{enumerate}
\end{finding}


%=============================================================================
\section{P2+P11 Hybrid: Sensitivity with Material-Optimised MEMS}
\label{sec:p2p11_materials}
%=============================================================================

From the W3i5 P10/P11 hybrid formula:
\begin{equation}
  \lam_{\rm hybrid} = \sqrt{
  \frac{2\muG c^4\,\Meff^{\rm MEMS}\,\Opsi\,\hbar}
  {G_{00}^2\,U_{\rm src}^2\,m^2\,Q_{\rm mech}^2\,d^{-2}\,T_{\rm int}}
  }.
  \label{eq:hybrid_formula_ref}
\end{equation}

With material-optimised MEMS (Si$_3$N$_4$, Phase 2):
\begin{itemize}[nosep]
  \item $\Meff^{\rm MEMS} = 10^{-9}$~kg (1~ng Si$_3$N$_4$ membrane)
  \item $Q_{\rm mech} = 10^9$ (dissipation-diluted, 20~mK)
  \item Other parameters as W3i5: $G_{00} = 0.90$, $U_{\rm src} = 500$~J,
    $m = 0.5$, $d = 0.5$~m, $T_{\rm int} = 10^6$~s
\end{itemize}

The improvement from $Q_{\rm mech} = 10^6$ (W3i5 baseline) to
$Q_{\rm mech} = 10^9$ (Si$_3$N$_4$) gives:
\begin{equation}
  \lam_{\rm min}^{\rm Si_3N_4} = \lam_{\rm min}^{\rm W3i5}
  \times \frac{10^6}{10^9} = 1.5\ee{-12} \times 10^{-3}
  = 1.5\ee{-15}.
\end{equation}

\begin{finding}[Material-Optimised P2+P11 Hybrid]
\label{find:optimised_hybrid}
With Si$_3$N$_4$ dissipation-diluted MEMS ($Q_{\rm mech} = 10^9$,
$m = 1$~ng), the P2+P11 hybrid achieves:
\begin{equation}
  \boxed{\lam_{\rm min}^{\rm Si_3N_4} \approx 1.5\ee{-15},}
\end{equation}
which is $10^6\times$ below the SQMS bound and comparable to the
ESS Lund optimistic projection. This establishes the \textbf{Si$_3$N$_4$-based
P2+P11 hybrid as a competitive alternative to the 256-cavity passive
array}, using a single source cavity and a single MEMS detector.

The key enabling technology is the dissipation dilution mechanism in
high-stress Si$_3$N$_4$, which provides $10^3\times$ improvement in
$Q_{\rm mech}$ over conventional metallic MEMS at the same mass scale.
\end{finding}


%=============================================================================
\section{Updated Material Property Table}
\label{sec:property_table}
%=============================================================================

\begin{table}[H]
\centering
\caption{Comprehensive material properties relevant to DFD detection.
Properties are evaluated at $T_{\rm op} = 20$~mK unless noted.}
\label{tab:properties}
\renewcommand{\arraystretch}{1.2}
\begin{tabular}{p{2.2cm}rrrrrr}
\toprule
Property & Nb & Nb$_3$Sn & MgB$_2$ & YBCO & Si$_3$N$_4$ & Al \\
\midrule
$\Tc$ (K) & 9.25 & 18.3 & 39 & 93 & --- & 1.2 \\
$\Delta$ (meV) & 1.55 & 3.1 & 2.5/7 & $\sim$20 & --- & 0.17 \\
Gap symmetry & s & s & s+s & d & --- & s \\
$\lambdaL$ (nm) & 39 & 85 & 140 & 150 & --- & 50 \\
$\xi_{\rm sc}$ (nm) & 38 & 3.2 & 5.2 & 1.5 & --- & 1600 \\
$\kappa$ & 1.0 & 27 & 27 & 100 & --- & 0.03 \\
$Q_0$ (SRF) & $2\ee{11}$ & $2\ee{10}$ & $5\ee{9}$ & $10^7$ & --- & $10^{10}$ \\
$Q_{\rm mech}$ & $10^6$ & --- & --- & --- & $10^{10}$ & $10^7$ \\
$\alpha_{\rm th}$ (4~K) & $10^{-9}$ & $3\ee{-9}$ & $5\ee{-9}$ & $3\ee{-8}$ & $2\ee{-9}$ & $5\ee{-9}$ \\
\midrule
$\FOMactive$ & 3.2 & 92 & 125 & 965 & --- & --- \\
$\FOMpassive$ (rel.) & $10^{20}$ & $10^{19}$ & $10^{18}$ & $10^{13}$ & --- & $10^{19}$ \\
$\FOMMEMS$ (rel.) & $10^{7}$ & --- & --- & --- & $10^{18}$ & $10^{10}$ \\
\bottomrule
\end{tabular}
\end{table}


%=============================================================================
\section{Updated Materials Roadmap: Post-Paradigm-Shift Programme}
\label{sec:roadmap}
%=============================================================================

The W3i5 paradigm shift (passive superiority, $\Meff$ correction,
two-component $\psi$) fundamentally restructures the DFD materials
programme into three parallel tracks.

\subsection{Track 1: Passive Cavity Programme (Nb-based, HIGHEST PRIORITY)}

\textbf{Objective}: Tighten the passive bound on $\lam$ from the
current $\lambdabound$ to $<10^{-12}$.

\textbf{Material}: N-doped Nb (single cavities and arrays).

\begin{center}
\renewcommand{\arraystretch}{1.3}
\begin{tabular}{p{2cm}p{3cm}p{2.5cm}p{3cm}}
\toprule
Phase & Configuration & Projected $\lam$ & Material milestone \\
\midrule
1A (now) & Single SQMS Nb cavity & $1.56\ee{-9}$ & Achieved \\
1B (2026) & Improved Q stability & $5\ee{-10}$ & 400$^\circ$C bake protocol \\
2A (2027) & Dual-cavity cross-corr. & $10^{-11}$ & Matched Nb pair \\
2B (2028) & 16-cavity array & $10^{-10}$ (incoh.) & Reproducible N-doping \\
3 (2029) & 256-cavity coherent & $6\ee{-12}$ & Frequency matching \\
4 (2030+) & ESS Lund class & $10^{-13}$--$10^{-14}$ & Cross-correlation \\
\bottomrule
\end{tabular}
\end{center}

\textbf{Key materials R\&D}:
\begin{itemize}[nosep]
  \item Develop 300--400$^\circ$C vacuum baking protocol to eliminate TLS
    noise from Nb$_2$O$_5$ surface oxide.
  \item Characterise cavity-to-cavity Q-factor reproducibility for
    N-doped Nb from the same ingot (needed for cross-correlation).
  \item Demonstrate $Q_0 > 10^{12}$ at $T < 20$~mK (single cavity).
  \item Develop in-situ frequency tuning to $\delta f/f < 10^{-12}$ for
    coherent array operation.
\end{itemize}

\subsection{Track 2: MEMS Proof Mass Programme (Si$_3$N$_4$, HIGH PRIORITY)}

\textbf{Objective}: Demonstrate the P2+P11 hybrid with progressively
lower proof mass and higher $Q_{\rm mech}$.

\textbf{Materials}: Al (Phase 1), Si$_3$N$_4$ + Al metallisation (Phase 2+).

\begin{center}
\renewcommand{\arraystretch}{1.3}
\begin{tabular}{p{2cm}p{3cm}p{2.5cm}p{3cm}}
\toprule
Phase & MEMS type & Projected $\lam$ & Material milestone \\
\midrule
1 (2027) & Al drumhead, $\mu$g & $5\ee{-11}$ & Al on Si, SQUID \\
2 (2029) & Si$_3$N$_4$ membrane, ng & $1.5\ee{-15}$ & Dissipation dilution \\
3 (2031+) & PhC nanobeam, pg & $\sim 10^{-16}$ & Nanofabrication \\
\bottomrule
\end{tabular}
\end{center}

\textbf{Key materials R\&D}:
\begin{itemize}[nosep]
  \item Develop Si$_3$N$_4$ membrane fabrication compatible with SRF
    cryogenic environment (vibration isolation, magnetic shielding).
  \item Demonstrate $Q_{\rm mech} > 10^9$ at $T = 20$~mK for metallised
    Si$_3$N$_4$ (Al metallisation must not degrade the $Q$).
  \item Design SQUID coupling geometry for membrane mode readout
    (flux-coupled or capacitive transduction).
  \item Investigate phononic crystal soft-clamping designs that
    maximise $Q_{\rm mech}/m$ at GHz mode frequencies relevant to
    $\Opsi = 8.17\ee{9}$~rad/s.
\end{itemize}

\subsection{Track 3: Exotic Materials Programme (YBCO/Nb$_3$Sn, LOWER PRIORITY)}

\textbf{Objective}: Maintain readiness for active detection recovery
if future theoretical developments reduce the detection gap (e.g.,
resolution of the $\Qpsi$ question, discovery of a non-$\Meff$-limited
coupling channel).

\textbf{Materials}: YBCO, Nb$_3$Sn laminates, nickelates.

\begin{center}
\renewcommand{\arraystretch}{1.3}
\begin{tabular}{p{2cm}p{3cm}p{3cm}p{3cm}}
\toprule
Phase & Activity & Material & Justification \\
\midrule
Ongoing & Nb$_3$Sn coating R\&D & Nb$_3$Sn on Nb & Passive Rank 2 \\
Monitor & YBCO grain boundary & YBCO bicrystal & Active Rank 1 (future) \\
Monitor & Nickelate thin films & La$_3$Ni$_2$O$_7$ & Highest $\delta T_c$ \\
Monitor & MATBG devices & MATBG on hBN & Flat-band DFD test \\
\bottomrule
\end{tabular}
\end{center}

\textbf{Rationale for lower priority}: The $10^{12}$--$10^{17}$ active
detection gap means that even the best active material (YBCO,
$\FOMactive = 965$) cannot close the gap through material optimisation
alone. Active detection recovery requires a theoretical breakthrough
(anomalously high $\Qpsi$, new coupling channel, or DFD field equation
modification). Until such a breakthrough occurs, the exotic materials
programme is in monitoring/maintenance mode.


%=============================================================================
\section{Cross-Patent Material Implications}
\label{sec:cross_patent}
%=============================================================================

\subsection{P1 (bubble timing): material-independent}

The P1 bubble timing experiment measures nonreciprocal propagation
in an EM-free geometry. The material requirements are thermal stability
and timing precision, not superconductor properties. P1 remains the
highest-priority near-term DFD test and is unaffected by the materials
paradigm shift.

\subsection{P5/P13 (nonreciprocal timing): moderate material sensitivity}

Nonreciprocal timing requires stable optical or microwave paths. The
material requirements are dimensional stability (low CTE) and optical
surface quality. Fused silica and ultra-low-expansion glass (ULE) are
optimal; superconductor properties are irrelevant.

\subsection{P3 (quantum computing): coherence preservation}

P3 requires superconducting qubit materials with minimal added
decoherence from the psi-field. W3i4 showed that psi-induced decoherence
is $\sim 0.1\%$ of the qubit linewidth at $\lam = \lambdabound$,
confirming negligible impact. The P3 material choice (Al or Ta for
transmon qubits) is unchanged.

\subsection{P11 (SRF source cavity): Nb remains optimal}

The P11 source cavity provides the EM energy ($U_{\rm src} = 500$~J)
that drives the psi-mode. The key requirement is high stored energy,
which favours Nb for its demonstrated performance in high-gradient
TESLA cavities. Nb$_3$Sn is an alternative for 4.2~K operation.


%=============================================================================
\section{Ranked Findings}
\label{sec:ranked}
%=============================================================================

\begin{enumerate}
  \item \textbf{[RANK 1 --- NEW] Passive FOM Ranking Inverts Active Ranking.}
    The passive material FOM is dominated by $Q_0$, not $\kappa$. N-doped
    Nb ($\FOMpassive \sim 10^{20}$) is optimal; YBCO ($\FOMpassive \sim 10^{13}$)
    is the worst candidate. This is the \emph{opposite} of the active ranking.
    (Theorem~\ref{thm:passive_inversion})

  \item \textbf{[RANK 2 --- NEW] Si$_3$N$_4$ Optimal for MEMS Proof Mass.}
    High-stress silicon nitride with dissipation dilution achieves
    $Q_{\rm mech} > 10^9$ at nanogram mass scales, giving $\FOMMEMS \sim 10^{18}$,
    far exceeding metallic alternatives. The Si$_3$N$_4$-based P2+P11 hybrid
    projects to $\lam \sim 1.5\ee{-15}$. (Theorem~\ref{thm:mems_ranking},
    Finding~\ref{find:optimised_hybrid})

  \item \textbf{[RANK 3 --- CONFIRMED] Active FOM Unchanged by Two-Component.}
    The two-component decomposition $\psi = \psigrav + \psiEM$ does not
    alter the active FOM formula. The material ranking (YBCO, MgB$_2$,
    Nb$_3$Sn, Nb) is preserved, but its practical relevance is diminished
    by the $10^{12}$--$10^{17}$ active detection gap.
    (Finding~\ref{find:active_fom_unchanged})

  \item \textbf{[RANK 4 --- NEW] Three-Track Materials Roadmap.}
    The post-paradigm-shift programme splits into: (1) Passive/Nb
    (highest priority, near-term), (2) MEMS/Si$_3$N$_4$ (high priority,
    medium-term), (3) Exotic/YBCO (monitoring mode, pending theoretical
    breakthrough). (Section~\ref{sec:roadmap})

  \item \textbf{[RANK 5 --- NEW] YBCO Demoted for Near-Term Programme.}
    YBCO's d-wave gap nodes make it fundamentally unsuitable for passive
    Q-factor measurements. Its role is restricted to the long-term active
    recovery track. (Finding~\ref{find:ybco_demoted})

  \item \textbf{[RANK 6 --- NEW] Material-Optimised Hybrid Competitive with
    256-Cavity Array.}
    The Si$_3$N$_4$ MEMS + P11 source hybrid ($\lam \sim 1.5\ee{-15}$)
    achieves sensitivity comparable to the ESS Lund optimistic projection,
    using a single cavity rather than a 256-cavity array. This makes
    Track 2 a cost-effective alternative to Track 1 Phase 3+.
    (Finding~\ref{find:optimised_hybrid})
\end{enumerate}


%=============================================================================
\section{Open Items for W3i7}
\label{sec:open}
%=============================================================================

\begin{enumerate}[nosep]
  \item \textbf{Nb TLS characterisation at 20 mK}: Quantify the TLS noise
    floor ($\sigma_{1/Q}$) for N-doped Nb at dilution refrigerator
    temperatures. This sets the fundamental limit on Track 1.
  \item \textbf{Si$_3$N$_4$ metallisation $Q$ degradation}: Measure
    $Q_{\rm mech}$ of a 20~nm Al-metallised Si$_3$N$_4$ membrane at
    20~mK. The Al layer must not introduce additional loss that degrades
    the dissipation dilution advantage.
  \item \textbf{SQUID coupling geometry for membrane modes}: Design the
    flux transformer that couples Si$_3$N$_4$ membrane displacement to
    a SQUID input coil. Estimate the displacement sensitivity
    ($\text{m}/\sqrt{\text{Hz}}$) and compare to the P2+P11 signal.
  \item \textbf{Frequency matching for Nb array}: Determine the
    achievable cavity-to-cavity frequency reproducibility for N-doped Nb
    TESLA cavities from the same ingot. This sets the feasibility
    threshold for the 256-cavity coherent array.
  \item \textbf{$\Qpsi$ material dependence}: Does the psi-mode quality
    factor depend on the cavity wall material? If $\Qpsi$ varies between
    Nb and Nb$_3$Sn, this would provide a new material selection criterion
    for both passive and active approaches.
\end{enumerate}


%=============================================================================
\section{Summary}
\label{sec:summary}
%=============================================================================

\begin{tcolorbox}[colback=green!5!white, colframe=green!50!black,
  title=\textbf{W3i6 Materials Science: Key Conclusions}]
\begin{enumerate}[nosep]
  \item The active FOM ($\kappa\cdot\Tc^{1/2}$) is \textbf{unchanged}
    under two-component decomposition. YBCO still leads for active mode.
    But active detection has a $10^{12}$+ gap, so this ranking is
    currently academic.
  \item The \textbf{Passive Material FOM} is derived:
    $\FOMpassive = Q_0 \cdot S_{\rm dim} \cdot S_{\rm freq}$.
    N-doped Nb dominates ($\FOMpassive \sim 10^{20}$); YBCO is worst
    ($\sim 10^{13}$). \textbf{Passive ranking inverts the active ranking.}
  \item The \textbf{MEMS proof mass FOM} identifies high-stress Si$_3$N$_4$
    as optimal ($\FOMMEMS \sim 10^{18}$), enabled by dissipation dilution
    achieving $Q_{\rm mech} > 10^9$ at nanogram masses.
  \item The Si$_3$N$_4$-based P2+P11 hybrid achieves $\lam \sim 1.5\ee{-15}$,
    competitive with the 256-cavity passive array.
  \item Three-track roadmap: \textbf{Passive/Nb} (highest priority),
    \textbf{MEMS/Si$_3$N$_4$} (high priority), \textbf{Exotic/YBCO}
    (monitoring mode).
\end{enumerate}
\end{tcolorbox}

\end{document}
